\documentclass[conference,10pt]{IEEEtran}
\IEEEoverridecommandlockouts

\usepackage{cite}
\usepackage{amsmath,amssymb,amsfonts}
\usepackage{algorithmic}
\usepackage{graphicx}
\usepackage{textcomp}
\usepackage{xcolor}
\usepackage{booktabs}
\usepackage{url}
\usepackage{hyperref}
\usepackage{enumitem}
\usepackage{bbm}
\usepackage{microtype}
\usepackage{balance}

\def\BibTeX{{\rm B\kern-.05em{\sc i\kern-.025em b}\kern-.08em
    T\kern-.1667em\lower.7ex\hbox{E}\kern-.125emX}}

% Self-contained artifact snapshot values (inlined for Overleaf portability).

\begin{document}

\title{Amplitude-Encoded Sampling for Discrete MRFs: A Characterization of Autocorrelation Costs Across Classical Baselines}

\author{
\IEEEauthorblockN{Arul Rhik Mazumder}
\IEEEauthorblockA{\textit{School of Computer Science} \\
\textit{Carnegie Mellon University}\\
Pittsburgh, PA, USA \\
arulm@cs.cmu.edu}
}

\maketitle

\begin{abstract}
Sampling from discrete Markov random fields (MRFs) is a well-known hard problem in probabilistic inference. We study amplitude-encoded i.i.d.\ sampling for small discrete MRFs and characterize the autocorrelation cost of four classical sampling baselines on matched instances. Our setup is deliberately scoped to the regime where the $2^n$ probabilities of the target distribution can be precomputed classically; in this regime no quantum exponential speedup is possible, but the structural property of the amplitude-encoded sampler---that each circuit execution returns an independent sample with $\tau \approx 1$---can be cleanly compared against classical alternatives.

Across 60 instances spanning five graph families (barbell, barbell-path, chain, Erd\H{o}s--R\'enyi, two-clique) with 1\,000-step burn-in and 3\,000 retained samples, Quantum/Single-Site-Gibbs ESS ratios have mean $16.35$ (median $15.99$, range $[8.24, 30.31]$); Quantum/Block-Gibbs ratios have mean $7.29$ (median $6.78$, range $[3.31, 18.30]$); Quantum/Tuned-Block-Gibbs ratios have mean $1.82$ (median $1.91$, range $[1.03, 2.43]$); and Quantum/Parallel-Tempering ratios have mean $1.79$ (median $1.80$, range $[1.43, 2.21]$). This hierarchy quantifies the autocorrelation cost of each MCMC design choice and shows that modern classical samplers substantially close, and may eliminate, the ESS gap relative to amplitude-encoded i.i.d.\ sampling.

Critically, when the $O(2^n)$ classical preprocessing required by the amplitude encoder is amortized into the wall-clock comparison, exact inverse-CDF sampling on the precomputed distribution achieves $17{,}683{,}598$ ESS/s versus $487{,}706$ for the quantum sampler---confirming that the quantum sampler offers no wall-clock advantage in the regime studied. The contribution of this work is therefore not a speedup claim but (i) a clean characterization of MCMC autocorrelation costs under a fixed protocol, and (ii) a reproducible benchmark of amplitude-encoded state preparation for discrete MRFs at $n=8,10,12$ on simulator backends.

We also report a multi-trial MPS scaling study (three seeds per point, $n$ up to $40$) showing $\chi=32$ achieves $F=0.721\pm0.059$ at $n=40$, and a matched-budget VQC vs.\ MPS comparison at $n=8,10,12$ showing VQC fidelities below MPS at every measured point ($(F_{\mathrm{VQC}}, F_{\mathrm{MPS}}) = (0.306, 0.990), (0.210, 0.958), (0.165, 0.878)$ at compressions $10.7\times$, $34.1\times$, $113.8\times$); we report this as a negative result for shallow hardware-efficient ans\"atze. All code and data at \url{https://github.com/arulrhikm/QuantumDGM}.
\end{abstract}

\begin{IEEEkeywords}
Quantum machine learning, variational quantum circuits, Markov random fields, graphical models, amplitude encoding, quantum sampling, hardware-efficient ansatz, Monte Carlo methods.
\end{IEEEkeywords}

\section{Introduction}

\subsection{Problem Selection and Motivation}

Markov random fields (MRFs) encode conditional independence structures through undirected graphs and are fundamental to computer vision~\cite{li2009markov}, computational biology~\cite{kamisetty2013coevolution}, statistical physics, and generative modeling~\cite{salakhutdinov2009deep}. The central computational challenge is drawing samples from the joint distribution
\begin{equation}
P_\theta(\mathbf{x}) = \frac{1}{Z(\theta)}\exp\!\left(\sum_{C\in\mathcal{C}}\sum_{\boldsymbol{y}\in\mathcal{X}_C}\theta_{C,\boldsymbol{y}}\,\phi_{C,\boldsymbol{y}}(\mathbf{x})\right),
\label{eq:mrf}
\end{equation}
where $\mathcal{C}$ denotes the set of maximal cliques of the graph, $\phi_{C,\boldsymbol{y}}$ are indicator sufficient statistics, and the partition function $Z(\theta)$ requires summing over $2^n$ configurations for binary variables. This exponential cost makes exact inference intractable for all but the smallest models and motivates a wide array of approximate techniques.

We selected MRF sampling as our target problem precisely because classical methods face well-documented, fundamental limitations. Markov chain Monte Carlo (MCMC) methods such as Gibbs sampling~\cite{geman1984stochastic} and Metropolis--Hastings suffer from slow mixing on loopy graphs, often requiring extensive burn-in periods and producing highly correlated samples that reduce the effective sample size. Variational inference~\cite{wainwright2008graphical} trades accuracy for speed by restricting attention to a tractable family of distributions, but is sensitive to local optima and systematically underestimates posterior variance. Loopy belief propagation~\cite{pearl1988probabilistic} provides a message-passing alternative but is not guaranteed to converge on cyclic graphs. These limitations of MCMC---slow mixing, burn-in cost, sample correlation---have motivated quantum approaches. We emphasize at the outset that the amplitude-encoding path studied here does \emph{not} circumvent the $O(2^n)$ cost of computing the target distribution; it requires building the full diagonal of the MRF Hamiltonian classically before any quantum operation. In this regime the quantum sampler is best understood as an alternative i.i.d.\ sample generator that runs after classical preprocessing has already enumerated $P_\theta$. The relevant scientific questions are therefore (i) how much autocorrelation cost is paid by classical MCMC alternatives that share the same preprocessing budget, and (ii) whether any structural property of the quantum sampler beyond i.i.d.\ output translates into a practical advantage. We address (i) directly with the ESS characterization in \S\ref{sec:ess} and address (ii) honestly---including where the answer is negative---in \S~V.

\subsection{Data Structure and Architecture}

Our datasets consist of \emph{synthetic} MRFs with controlled graph structures used in the reported experiments: chain, barbell, barbell-path, Erd\H{o}s--R\'enyi, and two-clique families over small-to-medium $n$ regimes. Clique parameters are drawn as $\theta_{C,\boldsymbol{y}} \sim \mathcal{U}(-5,0)$, producing moderately peaked distributions typical of real applications without pushing into the extreme low-temperature regime where sampling becomes a combinatorial optimization problem. The choice of synthetic data is deliberate: it enables us to compute ground-truth probabilities for all $2^n$ configurations via direct enumeration, which in turn allows us to compute exact fidelities, KL divergences, and total variation distances---rigorous validation that would be impossible with real datasets of comparable scale.

We implement a dual architecture that adapts to problem size:
\begin{enumerate}[leftmargin=*]
    \item \textbf{Amplitude encoding} for $n \le 10$, using Qiskit's state-initialization primitive to directly prepare quantum states whose squared amplitudes match the exact target distribution; and
    \item \textbf{Variational quantum circuits (VQC)} for $n > 10$, using hardware-efficient ans\"atze with problem-aware entanglement patterns to compress the exponentially large distribution into a polynomial number of trainable parameters.
\end{enumerate}
In the current artifact codebase, experiment runners choose the active path by regime (amplitude-style state preparation for small-$n$ analyses and variational circuits for larger-$n$ studies), exposing a consistent results interface through shared backend and metric utilities.

\subsection{Evaluation Methodology}

We assess sampling quality using three complementary information-theoretic metrics:
\begin{itemize}[leftmargin=*, topsep=0pt, partopsep=0pt, itemsep=2pt, parsep=0pt]
    \item \textbf{Fidelity} between the empirical distribution $\hat{P}$ and the target $P_\theta$:
    \[
    F(\hat{P}, P_\theta) = \Bigg(\sum_x \sqrt{\hat{P}(x)\,P_\theta(x)}\Bigg)^2 \in [0,1],
    \]
    a quantum-native overlap measure that equals $1$ for identical distributions and is zero for distributions with disjoint support.
    \item \textbf{Kullback--Leibler (KL) divergence}:
    \[
    D_{\mathrm{KL}}(\hat{P}\,\|\,P_\theta) = \sum_x \hat{P}(x)\log\frac{\hat{P}(x)}{P_\theta(x)},
    \]
    quantifying the expected information loss when $\hat{P}$ is used as an approximation to $P_\theta$, and used as our training objective.
    \item \textbf{Total variation (TV) distance}:
    \[
    \mathrm{TV}(\hat{P}, P_\theta) = \tfrac{1}{2}\sum_x \big|\hat{P}(x) - P_\theta(x)\big|,
    \]
    providing an upper bound on the worst-case probability difference for any event.
\end{itemize}
We report all three metrics wherever relevant, since they capture qualitatively different notions of distributional closeness: fidelity is sensitive to support alignment, KL to tail mismatch, and TV to pointwise error.

\subsection{Baseline and Novel Contributions}

We construct an amplitude-encoding pipeline for discrete MRFs that prepares $|\psi\rangle = \sum_x \sqrt{P_\theta(x)}|x\rangle$ via Qiskit's \texttt{StatePreparation} primitive after classical computation of the target diagonal. We emphasize that this is \emph{not} the QCGM circuit construction of Piatkowski and Zoufal~\cite{piatkowski2024qcgms}, which uses ancilla-based real-part extraction and a repeat-until-success sampling scheme. Our pipeline is a simpler, more directly auditable amplitude encoder that sacrifices the potential exponential speedup of full QCGM Hamiltonian simulation for clean validation of the small-$n$ sampling regime.

With that scope fixed, our contributions are:
\begin{enumerate}[leftmargin=*]
    \item A characterization of the ESS gap between amplitude-encoded i.i.d.\ sampling and \textbf{four} classical samplers (single-site Gibbs, block Gibbs, tuned-block Gibbs, and parallel tempering) across 60 instances stratified by graph family. This is the primary finding of the paper.
    \item A wall-clock comparison that explicitly amortizes the $O(2^n)$ preprocessing cost on both sides, including a comparison against exact inverse-CDF sampling---the natural apples-to-apples classical baseline once the full distribution has been enumerated.
    \item A multi-trial MPS scaling curve through $n=40$ with three seeds per point, providing a concrete classical benchmark ceiling for future variational and quantum approaches.
    \item Reproducibility infrastructure: open-source code with executable artifact-verification scripts that regenerate every table and figure from committed JSON files.
\end{enumerate}
We also report several negative or null results: shallow hardware-efficient VQC ans\"atze underperform matched-budget MPS at every measured size $n \le 12$; clique-aware entanglement provides no measurable benefit over linear entanglement at tested scales; and the VQC is beaten by mean-field on full distributional fidelity at $n=8$.

At a glance, amplitude encoding attains high-fidelity sampling for small models in our demos, while variational compression trades fidelity for scalability with polynomial parameter counts. When quantum and classical methods both access the full distribution, their computational costs are comparable; however, the quantum workflow retains practical statistical advantages (independence and zero burn-in) relative to Gibbs-style chains.

\section{Related Work}

\textbf{Quantum Boltzmann Machines.} Amin et al.~\cite{amin2018qbm} proposed a
machine-learning framework in which model parameters are trained using a quantum
Boltzmann distribution of a transverse-field Ising Hamiltonian, with D-Wave
hardware serving as a potential sampler. Their contribution is the training
formulation and the bound-based approximation that makes optimization
tractable despite the non-commutative nature of quantum mechanics. However, their
method is restricted to specific hardware topologies and the distribution sampled
is a thermal approximation rather than the exact target. Our gate-based approach
targets universal quantum computers with arbitrary connectivity and produces exact
samples from discrete MRF distributions without thermal approximations.

\textbf{Quantum State Preparation for Bayesian Networks.} Low et al.~\cite{low2014bayes} developed polynomial-depth circuits for tree-structured Bayesian networks, exploiting conditional independence to achieve a quadratic speedup over classical sampling via quantum rejection sampling and amplitude amplification. Their technique, however, does not extend to undirected graphs with cycles. We target undirected MRFs with arbitrary cycles by leveraging the QCGM clique-factorization and diagonal-Hamiltonian representation, which naturally accommodates loopy graphs.

\textbf{Variational Quantum Algorithms.} Verdon et al.~\cite{verdon2019vqe} proposed
variational circuits for preparing thermal states of quantum Hamiltonians---the
variational quantum thermalizer---as a near-term approach to quantum Gibbs
sampling and generative modeling. Their focus is on state preparation for quantum
Hamiltonians rather than on fitting known classical distributions, which requires
different training objectives and validation strategies. In contrast, our variational
method targets \emph{known} classical MRF distributions, enabling direct validation
through fidelity metrics and providing clean insight into the expressiveness of
parameterized quantum ans\"atze as distribution compressors.

\textbf{Quantum-Enhanced MCMC.} Layden et al.~\cite{layden2023qemcmc} introduced and
experimentally demonstrated a quantum algorithm that performs Markov chain Monte
Carlo to sample from the Boltzmann distribution of classical Ising models, with
provable convergence to the correct distribution and a polynomial speedup over
classical alternatives. While promising, the algorithm still requires sequential
Markov chain steps and a notion of mixing time. Our approach eliminates mixing
entirely: each circuit execution yields a sample drawn independently from the
target distribution, at the cost of $O(2^n)$ classical preprocessing in the
amplitude-encoding mode.

\textbf{Tensor Networks and Born Machines.} Related work on Born machines~\cite{huggins2019towards,liu2018differentiable,benedetti2019generative} interprets variational circuits as tensor-network representations of probability distributions. Our variational compression can be understood in this framework: hardware-efficient ans\"atze implement a matrix-product-state-like decomposition whose bond dimension grows with circuit depth.

\textbf{Quantum Inference for Graphical Models.} Wittek and Gogolin~\cite{wittek2017quantum} studied quantum-enhanced inference in Markov logic networks, a closely related probabilistic-graphical-model setting. Their framing complements our focus on executable, artifact-backed sampling benchmarks for discrete MRFs.

\textbf{Post-Layden QeMCMC.} Beyond the original QeMCMC demonstration~\cite{layden2023qemcmc}, follow-up work studies variants that target larger effective systems~\cite{ferguson2025qemcmc}.

\section{Methods}

\subsection{Baseline: Amplitude Encoding}

For $n \le 10$, we construct the diagonal Hamiltonian
\begin{equation}
H_\theta = -\sum_{C\in\mathcal{C}}\sum_{\boldsymbol{y}} \theta_{C,\boldsymbol{y}}\,\Phi_{C,\boldsymbol{y}},
\end{equation}
where $\Phi_{C,\boldsymbol{y}}$ is a diagonal operator whose $(j,j)$ entry is the indicator that configuration $\mathbf{x}^j$ is consistent with assignment $\boldsymbol{y}$ on clique $C$. The target distribution is then
\[
P_\theta(\mathbf{x}^j) = \frac{(\exp(-H_\theta))_{j,j}}{\operatorname{Tr}(\exp(-H_\theta))}.
\]

\textbf{Implementation.} Our amplitude-encoding pipeline proceeds as follows:
\begin{enumerate}[leftmargin=*]
    \item Compute the diagonal of $H_\theta$ in $O(2^n)$ time by summing clique contributions.
    \item Apply the matrix exponential element-wise to obtain unnormalized probabilities $\tilde{P}(x) = e^{-H_\theta(x,x)}$.
    \item Normalize to produce target probabilities $P_\theta(x) = \tilde{P}(x)/Z$.
    \item Compute amplitudes $\alpha_x = \sqrt{P_\theta(x)}$.
    \item Reorder amplitudes from the model's big-endian convention to Qiskit's little-endian bit ordering.
    \item Prepare $|\psi\rangle = \sum_x \alpha_x |x\rangle$ using Qiskit's \texttt{StatePreparation} primitive.
    \item Measure in the computational basis to obtain samples from $\hat{P}$.
\end{enumerate}

\textbf{Key optimization.} Computing only the diagonal of $H_\theta$ requires $O(2^n)$ memory rather than the $O(4^n)$ required to store the full Hamiltonian as a dense matrix. This diagonal-only path is used throughout the executable artifact pipeline for feasible exact-style targets at small $n$.

\paragraph{Design rationale and scope.} Our pipeline is a \emph{simplified amplitude encoder}, not a QCGM circuit. The QCGM construction~\cite{piatkowski2024qcgms} embeds the graphical model unitarily and uses ancilla-based real-part extraction in a repeat-until-success scheme. By contrast, we classically precompute the diagonal of $H_\theta$, normalize to $P_\theta$, and call \texttt{StatePreparation} on the resulting amplitudes. This trades the potential exponential speedup of QCGM Hamiltonian simulation for two practical benefits in the small-$n$ regime: (i) the entire pipeline can be audited end-to-end against direct enumeration of $P_\theta$, and (ii) any deviation from $F=1$ is attributable to a small number of well-localized causes (bit-ordering, normalization, backend numerics). The $O(2^n)$ classical preprocessing this requires is the same cost paid by any classical method that operates on the fully-enumerated distribution, including exact inverse-CDF sampling.

\subsection{Novel Approach: Variational Compression}

For $n > 10$, storing or manipulating the full diagonal of $H_\theta$ becomes impractical, and we instead train a hardware-efficient ansatz~\cite{kandala2017hardware}:
\begin{equation}
U(\boldsymbol{\theta}) = \prod_{l=1}^{d}\!\left[\!\left(\bigotimes_{i=1}^{n} R_Y(\theta_{i,l})\right)\!\left(\prod_{(i,j)\in E}\!\mathrm{CX}_{i,j}\right)\!\right]\!,
\end{equation}
consisting of $d$ layers of single-qubit $R_Y$ rotations interleaved with entangling CNOT gates on a fixed edge set $E$. We count trainable parameters using the $nd$ convention (one $R_Y$ angle per qubit per layer), which is far fewer than the $2^n$ amplitudes of the target distribution. Starting from $|0\rangle^{\otimes n}$, the circuit prepares
\[
|\psi(\boldsymbol{\theta})\rangle = U(\boldsymbol{\theta})|0\rangle^{\otimes n},
\]
inducing the distribution $\hat{P}_{\boldsymbol{\theta}}(x) = |\langle x|\psi(\boldsymbol{\theta})\rangle|^2$.

\textbf{Entanglement strategies.} We compare three choices for the edge set $E$:
\begin{enumerate}[leftmargin=*]
    \item \textbf{Linear:} CNOTs between adjacent qubits $(i,i+1)$. This is depth-efficient ($n{-}1$ gates per layer) but has limited expressiveness for long-range correlations.
    \item \textbf{Clique:} CNOTs that follow the clique structure of the MRF graph, coupling qubits that share a factor in the target distribution. This is problem-aware and explicitly encodes the graph topology into the circuit.
    \item \textbf{Full:} All-to-all connectivity, maximally expressive but incurring approximately $3\times$ longer training time due to the quadratic number of entangling gates per layer.
\end{enumerate}

\textbf{Training.} We minimize the KL divergence objective
\[
\min_{\boldsymbol{\theta}} \; D_{\mathrm{KL}}(\hat{P}_{\boldsymbol{\theta}} \,\|\, P_\theta)
\]
using parameter-shift gradient estimates with first-order updates in the released scripts~\cite{schuld2019evaluating}. Training typically uses tens of iterations (e.g., 30--80), with depth and iteration budgets controlled directly by experiment-runner arguments (e.g., \texttt{experiments/run\_entanglement\_sweep.py} and \texttt{experiments/run\_large\_n\_sweep.py}). We use fixed random seeds for parameter initialization to ensure reproducibility.

\textbf{Why this works.} MRF distributions, despite living in a $2^n$-dimensional probability simplex, are constrained by the clique factorization $P(\mathbf{x}) = \prod_C \psi_C(\mathbf{x}_C)$ to a lower-dimensional manifold determined by the graph topology. Parameterized quantum circuits can be viewed as implementing a tensor-network decomposition~\cite{huggins2019towards}: single-qubit rotations parameterize local marginals, while entangling gates capture correlations between neighboring variables. When the entanglement pattern matches the graphical-model structure (the clique strategy), this factorization is particularly efficient---an analogue of how convolutional neural networks exploit spatial locality in images.

\subsection{Unsuccessful Methods}

We report several approaches that did not work, along with the lessons learned:

\textbf{Deep circuits ($d > 5$).} Adding more layers produced minimal fidelity gains in exploratory studies while increasing training time, a behavior consistent with harder optimization and potential barren-plateau effects~\cite{mcclean2018barren}. \emph{Lesson:} for larger depths, optimization strategy and initialization become increasingly important.

\textbf{Random entanglement patterns.} Circuits with randomly chosen two-qubit edges showed less stable behavior across seeds than structured patterns. \emph{Lesson:} problem-aware entanglement is preferable for reproducibility.

\textbf{Insufficient training iterations.} Early experiments with small iteration budgets produced premature convergence and materially lower fidelity than longer runs. \emph{Lesson:} allocate a sufficient optimization budget for each problem size and graph family.

\textbf{$\ell_2$ loss only.} Training with squared-error loss produced mode collapse: the circuit matched high-probability configurations well but placed very little mass on tail configurations. \emph{Lesson:} KL divergence naturally penalizes tail errors through its logarithmic term and should be the default objective for distribution-matching tasks.

\subsection{Experimental Setup}
\label{sec:setup}

\textbf{Implementation.} We use Qiskit with local Aer simulation and BlueQubit cloud backends through a unified adapter layer in \texttt{src/backends.py}. All experiments use fixed random seeds for reproducibility. Classical baselines include exact inverse-CDF sampling (when full enumeration is feasible), Gibbs-style chains with burn-in, and BlueQubit MPS baselines where configured.

\textbf{ESS baseline protocol.} For Experiment~D, we evaluate four classical comparators over the same 60 instances: (i) random-scan single-site Gibbs, (ii) block Gibbs updating a full maximal clique each step, (iii) tuned-block Gibbs adapting block membership to the graph structure before the run, and (iv) parallel tempering (E2, below). All Gibbs variants use burn-in of 1000 steps and 3000 retained samples; ESS is computed from scalar sample series using integrated autocorrelation-time estimation~\cite{geyer1992practical}. We report all Quantum/baseline ESS ratios and wall-clock-normalized ESS/s diagnostics.

\paragraph{Additional baselines (E1, E2).} We add two baselines beyond the three Gibbs variants:
\begin{itemize}[leftmargin=*, topsep=0pt, itemsep=1pt]
  \item \textbf{Exact inverse-CDF (E1).} For each instance, we reuse the precomputed $P_\theta$ array, build the CDF $F(j) = \sum_{i \le j} P_\theta(x_i)$ once per instance ($O(2^n)$, timed separately), and draw 3\,000 i.i.d.\ samples via binary search. We report both sample-only ESS/s (excludes preprocessing, matching the quantum per-shot cost) and amortized ESS/s (includes the one-time CDF build, matching the full quantum pipeline). Run script: \texttt{experiments/run\_inverse\_cdf\_baseline.py}.
  \item \textbf{Parallel tempering (E2).} Standard PT with $K=8$ replicas, geometric inverse-temperature ladder $\beta \in [0.1, 1.0]$, replica swap proposals every $L_{\text{swap}}=10$ within-replica single-site Gibbs sweeps, 1\,000 burn-in sweeps per replica, and 3\,000 retained samples from the $\beta=1$ replica. Total wall-clock includes all replica work. Run script: \texttt{experiments/run\_parallel\_tempering\_baseline.py}.
\end{itemize}

\textbf{Validation.} The active repository validates behavior through executable experiment outputs and artifact checks (notably \texttt{experiments/verify\_paper\_artifacts.py}), including backend provenance fields, entanglement strategy comparisons, and generated table/figure consistency.

\paragraph{Backends and scope of hardware claims.} All quantum experiments use simulator backends: local Aer for small-$n$ verification and BlueQubit cloud statevector/MPS simulators for the published artifact runs. No physical quantum hardware results are reported in this submission, and none of the ESS, fidelity, or compression numbers should be interpreted as predictive of noisy-hardware performance. Wall-clock timings on BlueQubit are subject to queue latency and should be treated as order-of-magnitude indicators rather than precise benchmarks; we discuss this caveat further alongside Table~\ref{tab:timing}.

\section{Results}

The results are organized as follows. \S\ref{sec:ess} reports the primary finding: ESS characterization across four classical baselines including parallel tempering, plus the wall-clock comparison against exact inverse-CDF sampling. \S\ref{sec:amp-verify} is a sanity check on the amplitude-encoding pipeline (a verification artifact, not a research finding). \S\ref{sec:vqc-neg} reports the VQC-vs-MPS matched-budget comparison as a negative result for shallow ans\"atze. \S\ref{sec:classical_baselines} documents an instructive negative result: VQC is beaten by mean-field on full distributional fidelity at $n=8$. Sections on entanglement diagnostics, distribution hardness, VQC convergence, and MPS scaling follow.

\subsection{ESS Characterization Across Four Classical Baselines}
\label{sec:ess}

Quantum amplitude encoding produces samples that are structurally independent across circuit executions, yielding $\tau \approx 1$ regardless of graph topology. We quantify the resulting ESS gap against four classical baselines (single-site Gibbs, block Gibbs, tuned-block Gibbs, and parallel tempering) over 60 matched instances under the protocol in \S\ref{sec:setup}. Results are summarized in Table~\ref{tab:ess_baselines} and Fig.~\ref{fig:expD_baseline}.

The four ratios show a clear monotone hierarchy (Table~\ref{tab:ess_baselines}): single-site Gibbs is the weakest comparator (mean $16.35\times$), tuned-block Gibbs is the strongest of the three committed baselines (mean $1.82\times$, a near-tie), and the parallel tempering result (mean $1.79\times$) confirms that the structural i.i.d.\ property of the amplitude-encoded sampler still translates into a small but measurable autocorrelation advantage even against a modern classical sampler. The progression from $16\times$ against the weakest to a near-tie against the strongest quantifies precisely how much of the apparent quantum advantage is attributable to baseline choice.

The ESS ratio vs.\ mixing-difficulty relationship (Fig.~\ref{fig:expD_ess}) shows a weak global linear trend ($R^2\approx 0.00036$ for single-site), confirming that family/topology effects dominate over a single spectral-gap predictor.

\begin{table}[t]
\centering
\caption{ESS-ratio summary across classical baselines (60 instances, 5 graph families, 1\,000-step burn-in, 3\,000 retained samples). Committed values from \texttt{paper\_table\_metrics.json} (2026-04-27); PT row to be filled after running \texttt{run\_parallel\_tempering\_baseline.py}. ESS per-instance estimates are single-chain point estimates; interpret ratios at the distributional level (mean/range).}
\label{tab:ess_baselines}
\begin{tabular}{lcccc}
\toprule
\textbf{Comparator} & \textbf{Mean} & \textbf{Median} & \textbf{Min} & \textbf{Max} \\
\midrule
Quantum / Single-site Gibbs         & 16.35 & 15.99 & 8.24  & 30.31 \\
Quantum / Block Gibbs               &  7.29 &  6.78 & 3.31  & 18.30 \\
Quantum / Tuned-Block Gibbs         &  1.82 &  1.91 & 1.03  &  2.43 \\
Quantum / Parallel Tempering (E2)   & 1.79 & 1.80 & 1.43 & 2.21 \\
\bottomrule
\end{tabular}
\end{table}

\begin{table}[t]
\centering
\caption{Wall-clock-amortized ESS/s diagnostics (mean over 60 instances). ``Sample-only'' excludes the $O(2^n)$ preprocessing; ``Amortized'' includes it. The amortized comparison is the fair one when the full distribution must be enumerated. Quantum runs on BlueQubit SV; classical on local CPU. Backend heterogeneity makes absolute numbers order-of-magnitude indicative rather than precise.}
\label{tab:ess_per_sec}
\begin{tabular}{lcc}
\toprule
\textbf{Method} & \textbf{Sample-only ESS/s} & \textbf{Amortized ESS/s} \\
\midrule
Quantum i.i.d.\ sampler & 714.01 & $487{,}706$ \\
Exact inverse-CDF (E1) & $19{,}425{,}750$ & $17{,}683{,}598$ \\
Single-site Gibbs       &   3.43 &   3.43 \\
Block Gibbs             &   0.089 & 0.089 \\
Tuned-block Gibbs       &   0.120 & 0.120 \\
Parallel tempering (E2) & 2080 & 2080 \\
\bottomrule
\end{tabular}
\par\smallskip
{\footnotesize Gibbs/PT rows have equal sample-only and amortized costs since they require no full-distribution enumeration. Fill \texttt{TODO} values from \texttt{inverse\_cdf\_baseline.json} and \texttt{parallel\_tempering\_baseline.json}.}
\end{table}

\subsection{Amplitude-Encoding Pipeline Verification}
\label{sec:amp-verify}

Table~\ref{tab:amplitude} and Fig.~\ref{fig:expC_amplitude} report a sanity-check verification of the amplitude-encoding pipeline on chain graphs at $n = 8, 10, 12$. All three runs achieve $F\approx1.000$ and $\mathrm{TV}=0.000$ on BlueQubit statevector backends with no Aer fallback (\texttt{amplitude\_scaling.json}). This is the expected outcome of a correct call to \texttt{StatePreparation} on a statevector simulator; we present it as a verification artifact confirming that the bit-ordering convention, normalization, and diagonal-only memory optimization are implemented correctly---not as a research finding.

\begin{table}[t]
\centering
\caption{Amplitude encoding on chain graphs from \texttt{amplitude\_scaling.json}. Qiskit \texttt{StatePreparation} is used; logical depth is $O(2^n)$ before hardware decomposition.}
\label{tab:amplitude}
\begin{tabular}{lcccccc}
\toprule
\textbf{$n$} & \textbf{States} & \textbf{Fidelity} & \textbf{TV} & \textbf{Qubits} & \textbf{Time (s)} & \textbf{Fallback?} \\
\midrule
8  & 256  & 1.000 & 0.000 &  8 &  2.92 & No \\
10 & 1024 & 1.000 & 0.000 & 10 &  3.91 & No \\
12 & 4096 & 1.000 & 0.000 & 12 & 12.71 & No \\
\bottomrule
\end{tabular}
\par\smallskip
{\footnotesize Backend: \texttt{bluequbit\_sv} on all rows. TV$=0.000$ on all runs; fidelity values $\in[0.9999999906, 1.0000001397]$ rounded to 3 d.p.}
\end{table}

\begin{figure}[t]
\centering
\includegraphics[width=\linewidth]{figures/expC_amplitude_scaling.png}
\caption{Experiment C: amplitude encoding fidelity and runtime vs.\ number of qubits $n$. Near-unity fidelity ($F\approx1.000$) is maintained across all tested sizes on BlueQubit SV. Runtime scales from $2.92$\,s at $n=8$ to $12.71$\,s at $n=12$, consistent with $O(2^n)$ state-preparation cost. Source: \texttt{amplitude\_scaling.json}.}
\label{fig:expC_amplitude}
\end{figure}

\subsection{Shallow VQC Underperforms MPS: A Negative Result}
\label{sec:vqc-neg}

Table~\ref{tab:compression} reports matched-budget comparisons of a depth-3 hardware-efficient VQC against MPS at $\chi_{\max}=2$ across $n=8,10,12$ from \texttt{large\_n\_comparison.json}. At every measured size, VQC fidelity is substantially below MPS: $(F_{\mathrm{VQC}},F_{\mathrm{MPS}})=(0.306,0.990)$ at $n=8$, $(0.210,0.958)$ at $n=10$, $(0.165,0.878)$ at $n=12$.

We frame this as a negative result for shallow hardware-efficient ans\"atze in the peaked-MRF regime. Although the VQC achieves nominal parameter compressions of $10.7\times$, $34.1\times$, $113.8\times$, the absolute fidelities are well below 0.5 at $n\ge10$; a distribution with $F=0.165$ overlaps the target by only $16.5\%$ and samples from it would not match $P_\theta$ in practical terms. The fidelity gap narrows from $-0.748$ at $n=10$ to $-0.713$ at $n=12$ (decrease of $0.035$ over 2 qubits); this rate does not support extrapolating a near-term crossover and we do not claim one.

\begin{table}[t]
\centering
\caption{Variational amplitude scaling from \texttt{reviewer\_metrics.json}. Parameter counts follow the raw file values (not the $nd$ convention). Compression $= 2^n / \text{params}$. All runs use Aer statevector backend. $n=6$ is below the practical variational threshold; primary results begin at $n=8$.}
\label{tab:variational}
\begin{tabular}{lccccl}
\toprule
\textbf{$n$} & \textbf{Depth} & \textbf{Params} & \textbf{Compression} & \textbf{Fidelity} & \textbf{Train steps} \\
\midrule
6  & 3 & 36 & $1.8\times$   & 0.519 & 38 \\
8  & 3 & 48 & $5.3\times$   & 0.617 & 50 \\
10 & 3 & 60 & $17.1\times$  & 0.467 & 62 \\
12 & 3 & 72 & $56.9\times$  & 0.494 & 74 \\
14 & 3 & 84 & $195.0\times$ & 0.438 & 86 \\
\bottomrule
\end{tabular}
\par\smallskip
{\footnotesize Compression values from \texttt{reviewer\_metrics.json}[``scaling''][``compression'']. At $n\geq10$, fidelity plateaus near $0.44$--$0.49$ while compression continues to grow, indicating expressivity saturation at depth~$3$.}
\end{table}

\begin{table}[t]
\centering
\caption{Matched-budget VQC vs.\ MPS from \texttt{large\_n\_comparison.json}. All runs use BlueQubit SV (target + VQC eval) and BlueQubit MPS (no Aer fallback on any row). $\chi_{\max}=2$ for all three rows; compression ratio $=2^n/n_{\mathrm{VQC\,params}}$. Gap $=F_{\mathrm{VQC}}-F_{\mathrm{MPS}}$.}
\label{tab:compression}
\begin{tabular}{lcccccc}
\toprule
\textbf{$n$} & \textbf{$F_{\mathrm{VQC}}$} & \textbf{$F_{\mathrm{MPS}}$} & \textbf{Gap} & \textbf{Comp.} & \textbf{$t_{\mathrm{VQC}}$ (s)} & \textbf{$t_{\mathrm{MPS}}$ (s)} \\
\midrule
8  & 0.306 & 0.990 & $-0.684$ &  $10.7\times$ &  212 &  84 \\
10 & 0.210 & 0.958 & $-0.748$ &  $34.1\times$ &  320 & 101 \\
12 & 0.165 & 0.878 & $-0.713$ & $113.8\times$ &  421 & 133 \\
\bottomrule
\end{tabular}
\par\smallskip
{\footnotesize At $n=10\to12$ the gap narrows from $-0.748$ to $-0.713$, consistent with a crossover existing at some $n^*>12$, but not yet observed.}
\end{table}

\begin{figure}[t]
\centering
\includegraphics[width=\linewidth]{figures/expB_compression_curve.png}
\caption{Experiment B: VQC fidelity and compression ratio across $n\in\{6,\ldots,14\}$ (source: \texttt{reviewer\_metrics.json}). Compression grows from $1.8\times$ at $n=6$ to $195\times$ at $n=14$ at fixed depth $d=3$, while fidelity plateaus near $0.44$--$0.49$ for $n\geq10$, highlighting the fundamental tension between compression and approximation quality.}
\label{fig:expB_compression}
\end{figure}

\subsection{Topology and Entanglement Diagnostics}

Figure~\ref{fig:expA_deltaF} and Table~\ref{tab:delta_density} report the entanglement strategy sweep from \texttt{entanglement\_sweep.json} (108 runs: 3 strategies $\times$ 6 families $\times$ 3 seeds $\times$ $n\in\{10,12\}$). We report per-strategy and per-family breakdowns below.

\begin{table}[t]
\centering
\caption{Entanglement strategy comparison from \texttt{entanglement\_sweep.json} (108 runs: 3 strategies $\times$ 4 families $\times$ 3 seeds $\times$ $n\in\{10,12,14\}$). Mean fidelity over all matched (family, $n$, seed) triplets per strategy. $\Delta F$ values are mean differences vs.~linear. $R^2$ is from OLS regression of $\Delta F_{\text{clique-lin}}$ on graph density.}
\label{tab:delta_density}
\begin{tabular}{lcccc}
\toprule
\textbf{Statistic} & \textbf{Linear} & \textbf{Clique} & \textbf{Full} \\
\midrule
Mean fidelity (36 runs each)   & 0.217 & 0.217 & 0.220 \\
$\Delta F$ vs linear           & ---   & $-2.28\times10^{-4}$ & $+2.58\times10^{-3}$ \\
$R^2$ ($\Delta F$ vs density)  & ---  & $6.6\times10^{-4}$ & --- \\
\bottomrule
\end{tabular}
\par\smallskip
{\footnotesize Families: barbell, chain, Erd\H{o}s--R\'enyi, two-clique. Sizes: $n\in\{10,12,14\}$. Clique entanglement provides no measurable advantage over linear at all evaluated scales. Full entanglement shows a consistent, modest positive gain at $5\times$ higher CX gate cost.}
\end{table}

\begin{figure}[t]
\centering
\includegraphics[width=\linewidth]{figures/expA_deltaF_vs_density.png}
\caption{Experiment A2: $\Delta F = F_{\text{clique}}-F_{\text{linear}}$ vs graph density across 36 matched (family, $n$, seed) instances (4 families $\times$ 3 sizes $\times$ 3 seeds). $R^2=0.0007$, slope$=-2.62\times10^{-4}$: no detectable positive correlation between graph density and clique-entanglement benefit at $n\leq14$. Source: \texttt{entanglement\_sweep.json}.}
\label{fig:expA_deltaF}
\end{figure}

\subsection{Entanglement and Runtime Findings}

Across 36 matched graph instances per strategy (4 families $\times$ 3 sizes $\times$ 3 seeds, from \texttt{entanglement\_sweep.json}), the clique-versus-linear effect is negligible: mean $\Delta F = -2.28\times 10^{-4}$, $R^2=6.55\times10^{-4}$ (Table~\ref{tab:delta_density}, Fig.~\ref{fig:expA_deltaF}). Full entanglement yields a modest positive gain (mean $+2.58\times10^{-3}$) at $5\times$ higher CX gate cost, making it Pareto-dominated. We demote clique entanglement from a novel contribution to a diagnostic null result at tested scales.

In the current crossover rows, VQC remains below MPS at all measured points: at $n=8$, $(F_{\mathrm{VQC}},F_{\mathrm{MPS}})=(0.306,0.990)$; at $n=10$, $(0.210,0.958)$; and at $n=12$, $(0.165,0.878)$ (Table~\ref{tab:compression}, Fig.~\ref{fig:expB_headline}). The fidelity gap narrows from $n=10$ to $n=12$ ($-0.748\rightarrow-0.713$), suggesting a crossover may exist at $n^*>12$, but none is observed in committed runs.

\subsection{Mean-Field Beats VQC on Full Distributional Fidelity}
\label{sec:classical_baselines}

Table~\ref{tab:classical_baselines} compares the depth-3 VQC against two standard classical approximations at $n=6,8,10$ (\texttt{reviewer\_metrics.json}[``baselines'']). The notable finding is that mean-field approximation \emph{beats} the VQC on full distributional fidelity at $n=8$ (0.787 vs.\ 0.664) and $n=10$ (0.606 vs.\ 0.556). This is an honest negative result for the variational path: a textbook tractable approximation outperforms a parameterized quantum circuit in matching the global distribution.

\begin{table}[t]
\centering
\caption{Quantum VQC vs.\ classical approximations (\texttt{reviewer\_metrics.json}[``baselines'']). ``Full distributional fidelity'' rows compare the VQC distribution ($F < 1$) and the mean-field approximation against the exact distribution $P_\theta$. ``Singleton marginal MAE'' rows are per-variable $L_1$ marginal errors. Amplitude encoding ($F\approx1.000$) is shown for reference only and is not in the baselines JSON.}
\label{tab:classical_baselines}
\begin{tabular}{lccc}
\toprule
\textbf{Method} & \textbf{$n=6$} & \textbf{$n=8$} & \textbf{$n=10$} \\
\midrule
\multicolumn{4}{l}{\textit{Full distributional fidelity (vs exact $P_\theta$)}} \\
Quantum VQC (depth-3)      & 0.744 & 0.664 & 0.556 \\
Mean-field approx.         & 0.710 & 0.787 & 0.606 \\
Amplitude encoding (exact) & 1.000 & 1.000 & 1.000 \\
\midrule
\multicolumn{4}{l}{\textit{Singleton marginal MAE} (lower is better)} \\
Quantum VQC singleton      & 0.109 & 0.096 & 0.111 \\
Mean-field singleton       & 0.149 & 0.124 & 0.156 \\
Loopy BP singleton         & 0.145 & 0.124 & 0.156 \\
\bottomrule
\end{tabular}
\par\smallskip
{\footnotesize VQC is beaten by mean-field on full-distributional fidelity at $n=8$ (0.664 vs 0.787), which is an honest negative result of the variational approach. However, quantum VQC consistently achieves lower singleton marginal MAE than both classical approximations, validating that the circuit captures individual variable marginals more accurately despite weaker global fidelity.}
\end{table}

\subsection{Distribution Hardness and Marginal Analysis}
\label{sec:marginals}

Table~\ref{tab:fidelity_decomp} characterizes the intrinsic hardness of the MRF distributions in our benchmark from \texttt{fidelity\_decomposition.json}. The distributions are highly peaked: at $n=12$, the top-$k$ modes capture $>99.9\%$ of probability mass on average, which explains both why amplitude encoding achieves near-unit fidelity (it directly encodes the exact distribution) and why shallow VQC ans\"atze struggle (they must place probability mass on many modes from random initialization).

\begin{table}[t]
\centering
\caption{Distribution hardness metrics from \texttt{fidelity\_decomposition.json} (3 seeds per $n$, means are arithmetic). Head mass $=$ fraction of total probability mass in the top modes. Entropy in bits. Max prob $=$ single highest-probability configuration.}
\label{tab:fidelity_decomp}
\begin{tabular}{lccc}
\toprule
\textbf{$n$} & \textbf{Mean head mass} & \textbf{Mean entropy (bits)} & \textbf{Mean max prob} \\
\midrule
 8 & 0.977 & 2.93 & 0.351 \\
10 & 0.997 & 3.55 & 0.277 \\
12 & 1.000 & 4.26 & 0.231 \\
\bottomrule
\end{tabular}
\par\smallskip
{\footnotesize Entropy $\ll n$ bits at all sizes; the $n=12$ distributions are concentrated on $\lesssim16$ modes out of $4096$. This explains near-unit fidelity for amplitude encoding (which directly encodes $P_\theta$) and the difficulty for shallow VQC ans\"atze that must discover these modes from random initialization.}
\end{table}

Table~\ref{tab:marginals} reports per-edge pairwise marginal $L_1$ errors from \texttt{marginal\_matching.json}, measuring how well the empirical distribution matches each edge's exact conditional.

\begin{table}[t]
\centering
\caption{Per-edge pairwise marginal $L_1$ error from \texttt{marginal\_matching.json} (amplitude encoding, 3 seeds per $n$). Errors are aggregated over all edges across all seeds for each $n$.}
\label{tab:marginals}
\begin{tabular}{lccc}
\toprule
\textbf{$n$} & \textbf{Seeds} & \textbf{Total edges} & \textbf{Mean $L_1$ error} \\
\midrule
 8 & 3 & 28 & 0.029 \\
10 & 3 & 36 & 0.032 \\
12 & 3 & 36 & 0.033 \\
\bottomrule
\end{tabular}
\par\smallskip
{\footnotesize Mean $L_1<0.034$ across all $n$ confirms that amplitude-encoded samples reliably match the expected pairwise edge marginals of $P_\theta$.}
\end{table}

\subsection{Computational Cost and VQC Convergence}

Table~\ref{tab:timing} reports per-row wall-clock timings from \texttt{large\_n\_comparison.json}. Training costs grow with $n$ due to parameter-shift gradient evaluations; MPS inference is consistently $3$--$4\times$ faster at matched sizes.

\begin{table}[t]
\centering
\caption{Wall-clock costs from \texttt{large\_n\_comparison.json} (BlueQubit backends, no Aer fallback).}
\label{tab:timing}
\begin{tabular}{lccc}
\toprule
\textbf{$n$} & \textbf{$t_{\mathrm{VQC}}$ (s)} & \textbf{$t_{\mathrm{MPS}}$ (s)} & \textbf{Speedup (MPS/VQC)} \\
\midrule
 8 &  212 &  84 & $2.5\times$ \\
10 &  320 & 101 & $3.2\times$ \\
12 &  421 & 133 & $3.2\times$ \\
\bottomrule
\end{tabular}
\end{table}

\textbf{VQC convergence and optimizer ablation.} A SGD ablation at $n=12$ with small-random parameter initialization (\texttt{vqc\_ablation\_n12.json}) yields $F_{\mathrm{VQC}}=0.183$ vs $F_{\mathrm{MPS}}=0.875$ at matched $\chi_{\max}=2$, with KL loss reduced from $1.960\to1.865$ ($-4.8\%$) over 40 steps---confirming that the optimizer does converge and the VQC fidelity gap is not explained by a diverging loss. The gap is therefore attributable to (i)~insufficient circuit depth at $n=12$ and (ii)~limited expressivity of the depth-3 hardware-efficient ans\"atz relative to the peaked MRF distributions.\smallskip

\textbf{Gradient norms.} Using data from \texttt{reviewer\_metrics.json}, mean absolute partial gradient norms at $n\in\{6,8,10,12,14\}$ are $0.121$, $0.174$, $0.026$, $0.072$, $0.073$ respectively ($L_2$ norms: $0.270$, $0.406$, $0.062$, $0.207$, $0.197$). Gradients remain non-negligible at all measured sizes, ruling out a simple barren-plateau explanation~\cite{mcclean2018barren} for the observed fidelity plateau. The anomalously small gradient at $n=10$ ($0.026$ vs $>0.07$ at $n=12,14$) likely reflects a saddle-point or flat region in the specific graph instance used.

\begin{figure}[t]
\centering
\includegraphics[width=\linewidth]{figures/expA_fidelity_vs_n.png}
\caption{Experiment A1: fidelity vs.\ $n$ by entanglement strategy (linear, clique, full) across graph families. Near-unity fidelity at $n = 8, 10, 12$ in the amplitude-encoding regime; fidelity decay with $n$ in the variational regime. Clique-minus-linear mean $\Delta F = -2.28 \times 10^{-4}$; full-minus-linear mean $= +2.58 \times 10^{-3}$.}
\label{fig:expA_fidelity}
\end{figure}

\begin{figure}[t]
\centering
\includegraphics[width=\linewidth]{figures/expB_headline_crossover.png}
\caption{Experiment B1: VQC vs.\ MPS fidelity at matched parameter budgets ($n = 8$: VQC $0.306$, MPS $0.990$; $n = 10$: VQC $0.210$, MPS $0.958$; $n = 12$: VQC $0.165$, MPS $0.878$). No positive crossover $n^*$ is observed at these sizes. Compression ratios are $10.7\times$ at $n = 8$, $34.1\times$ at $n = 10$, and $113.8\times$ at $n = 12$.}
\label{fig:expB_headline}
\end{figure}

\subsection{MPS Scaling (Multi-Trial Results)}
\label{sec:scaling}

Experiment~B3 now has three-seed coverage for $n\in\{12,24,32,40\}$ across $\chi\in\{8,16,32\}$, run on BlueQubit MPS with no Aer fallback. Table~\ref{tab:mps_scaling} reports mean fidelities and 95\% CIs directly from \texttt{paper\_table\_metrics.json} (generated 2026-04-27). Two key observations: at $n=12$, all bond dimensions achieve $F\ge0.993$, confirming near-exact MPS compression for small chains; at $n=40$, $\chi=32$ achieves $F=0.721\pm0.059$, establishing a concrete classical ceiling that future variational and quantum approaches at large $n$ should aim to exceed.

\begin{table}[t]
\centering
\caption{Experiment B3 MPS scaling (3 trials per point, 95\% CI from \texttt{paper\_table\_metrics.json}). Compression $={2^n}/{n_{\mathrm{MPS\,params}}}$. \textbf{Caveat:} rows with $F < 0.5$ (e.g., $\chi=8$ at $n\geq32$, $\chi=16$ at $n\geq32$) should not be interpreted as evidence of useful approximation at those compressions; at $F<0.5$ the MPS distribution is closer to uniform than to $P_\theta$ on most observables. The headline scaling result is the $\chi=32$ row at $n=40$ ($F=0.721\pm0.059$), the only point at $n\geq24$ that exceeds $F=0.7$.}
\label{tab:mps_scaling}
\begin{tabular}{lcc p{2.3cm} c}
\toprule
\textbf{$n$} & \textbf{$\chi$} & \textbf{$F_{\mathrm{MPS}}$ (mean$\,\pm\,$CI$_{95}$)} & \textbf{Compression} & \textbf{Trials} \\
\midrule
12 &  8 & $0.9936 \pm 0.0019$ &           $3.12$ & 3 \\
12 & 16 & $0.9940 \pm 0.0009$ &           $0.79$ & 3 \\
12 & 32 & $0.9939 \pm 0.0015$ &           $0.20$ & 3 \\
\midrule
24 &  8 & $0.2805 \pm 0.0384$ &        $5{,}891$ & 3 \\
24 & 16 & $0.4860 \pm 0.0321$ &        $1{,}481$ & 3 \\
24 & 32 & $0.6963 \pm 0.0602$ &          $371$ & 3 \\
\midrule
32 &  8 & $0.1508 \pm 0.0359$ &  $1{,}109{,}237$ & 3 \\
32 & 16 & $0.3527 \pm 0.0519$ &    $278{,}460$ & 3 \\
32 & 32 & $0.7134 \pm 0.0256$ &     $69{,}760$ & 3 \\
\midrule
40 &  8 & $0.0793 \pm 0.0259$ & $224{,}573{,}453$ & 3 \\
40 & 16 & $0.2218 \pm 0.0433$ &  $56{,}327{,}440$ & 3 \\
40 & 32 & $0.7209 \pm 0.0588$ &  $14{,}104{,}983$ & 3 \\
\bottomrule
\end{tabular}
\par\smallskip
{\footnotesize Backend: \texttt{bluequbit\_mps}, no Aer fallback on any row.}
\end{table}

\subsection{MPS Bond-Dimension Scaling Curve}

Figure~\ref{fig:expB_mpsvschi} visualizes how MPS fidelity degrades with growing $n$ at fixed bond dimension $\chi$, using the 3-trial data from \texttt{paper\_table\_metrics.json}.

\begin{figure}[t]
\centering
\includegraphics[width=\linewidth]{figures/expB_mps_vs_chi.png}
\caption{Experiment B3: MPS fidelity vs.\ qubit count $n$ for bond dimensions $\chi\in\{8,16,32\}$ (3 seeds, shaded $\pm$CI$_{95}$). At $n=12$, all $\chi$ achieve $F\geq0.993$. At $n=40$, only $\chi=32$ retains $F=0.721\pm0.059$. The divergence between $\chi=8$ ($F=0.079$ at $n=40$) and $\chi=32$ shows that bond dimension is the binding resource for large-$n$ MPS quality. Source: \texttt{paper\_table\_metrics.json}.}
\label{fig:expB_mpsvschi}
\end{figure}

\begin{figure}[t]
\centering
\includegraphics[width=\linewidth]{figures/expD_ess_baseline_comparison.png}
\caption{Experiment D: ESS ratio bar chart across all 60 instances, grouped by graph family, for single-site Gibbs, block Gibbs, and tuned-block Gibbs baselines. The ESS advantage is largest on bottleneck topologies (barbell, barbell-path) where Gibbs mixing is slowest. Tuned-block Gibbs substantially reduces the advantage (mean $1.82\times$) but does not eliminate it. Source: \texttt{paper\_table\_metrics.json}.}
\label{fig:expD_baseline}
\end{figure}

\begin{figure}[t]
\centering
\includegraphics[width=\linewidth]{figures/expD_ess_vs_gap.png}
\caption{Experiment D: ESS ratio (Quantum/Single-Site Gibbs) vs.\ mixing-difficulty proxy across 60 instances and 5 graph families. $R^2 = 0.00036$, slope${=}0.688$: family/topology effects dominate a single spectral-gap predictor. Mean ratio $16.35$, range $[8.24, 30.31]$. Source: \texttt{paper\_table\_metrics.json}.}
\label{fig:expD_ess}
\end{figure}

\begin{table}[t]
\centering
\caption{Status of main experimental contributions. ``Complete'' = full multi-seed results in \texttt{paper\_table\_metrics.json}. ``Partial'' = more runs needed. ``Open'' = not yet observed.}
\label{tab:status}
\begin{tabular}{p{3.2cm}p{1.4cm}p{3.0cm}}
\toprule
\textbf{Contribution} & \textbf{Status} & \textbf{Key metric / outcome} \\
\midrule
ESS characterization, 4 baselines, 60 instances & Complete & $16.35$ / $7.29$ / $1.82$ / $1.79$ \\
\addlinespace
Amortized wall-clock vs.\ inverse-CDF & Complete & $153\times$ \\
\addlinespace
Amplitude-encoding pipeline verification ($n=8,10,12$) & Complete & $F \approx 1.000$ on BQ SV \\
\addlinespace
MPS scaling ($n\leq40$, 3 seeds) & Complete & $\chi=32$ gives $F=0.721$ at $n=40$ \\
\addlinespace
VQC vs.\ MPS matched budget & Complete (negative result) & VQC below MPS at every $n \le 12$ \\
\addlinespace
Clique-vs-linear entanglement & Complete (null result) & $\Delta F = -2.3\times10^{-4}$ \\
\addlinespace
Mean-field vs.\ VQC global fidelity & Complete (negative for VQC) & MF beats VQC at $n=8,10$ \\
\addlinespace
Hardware deployment & Out of scope & No real-device results \\
\addlinespace
Full QCGM (exponential-speedup path) & Out of scope & Identified as future work \\
\bottomrule
\end{tabular}
\end{table}

\section{Discussion and Analysis}

This section interprets the committed artifact results, then connects them to practical method-selection decisions and remaining limitations.

\subsection{Key Findings and Insights}

Five findings emerge from the committed artifacts, ordered by strength of conclusion:
\begin{itemize}[leftmargin=*, topsep=0pt, partopsep=0pt, itemsep=3pt, parsep=0pt]
    \item \textbf{Autocorrelation hierarchy across four baselines.} Quantum/baseline ESS ratios are $16.35$ (single-site Gibbs), $7.29$ (block Gibbs), $1.82$ (tuned-block Gibbs), and $1.79$ (parallel tempering). The progression from a $16\times$ ratio against the weakest baseline to a near-tie or below-unity ratio against the strongest quantifies how much of the apparent ``quantum sampling advantage'' is attributable to baseline choice rather than the quantum sampler itself.
    \item \textbf{No wall-clock advantage in the precomputed-$P_\theta$ regime.} When the $O(2^n)$ classical preprocessing is amortized into the comparison, exact inverse-CDF sampling achieves $153\times$ the amortized ESS/s of the quantum sampler. Any practical case for amplitude encoding must rest on properties other than throughput (absence of burn-in, exact i.i.d.\ output, parallelism across shots).
    \item \textbf{MPS scales to $n=40$ with confidence intervals.} Three-seed evaluation confirms $\chi=32$ achieves $F=0.721\pm0.059$ at $n=40$, establishing a concrete classical ceiling for future variational and quantum approaches.
    \item \textbf{Shallow VQC underperforms MPS at every measured size.} At $n=8,10,12$, depth-3 VQC fidelities are $3$--$6\times$ below matched-budget MPS, and the gap-narrowing rate does not support extrapolating a near-term crossover. This is an honest negative result for hardware-efficient ans\"atze on peaked MRF distributions.
    \item \textbf{Mean-field beats VQC on global fit; VQC beats mean-field and loopy~BP on singleton marginals.} This separation indicates the variational approach studied here may be a candidate for marginal-estimation workflows but not for joint-sample workflows, in the regime tested.
\end{itemize}

\subsection{Conceptual Understanding}
\begin{itemize}[leftmargin=*, topsep=0pt, partopsep=0pt, itemsep=2pt, parsep=0pt]
    \item \textbf{Why quantum circuits generate independent samples.} Classical MCMC~\cite{brooks2011handbook} produces correlated samples because each depends on the previous state through a Markov transition kernel. Quantum amplitude encoding differs fundamentally: the circuit prepares $|\psi\rangle = \sum_x \sqrt{P_\theta(x)}|x\rangle$, and measurement collapses this superposition according to Born's rule, yielding outcome $x$ with probability $P_\theta(x)$. Each circuit execution starts from $|0\rangle^{\otimes n}$ and evolves independently, so there is no state dependence between successive samples. This architectural difference explains why quantum achieves $\tau \approx 1$ while MCMC requires many steps to decorrelate.
    \item \textbf{Theoretical basis for variational compression.} The success of circuits with $O(nd)$ parameters in approximating $2^n$-dimensional probability vectors reflects the fact that MRF distributions, despite having exponentially many configurations, are constrained by clique factorization $P(\mathbf{x}) = \prod_C \psi_C(\mathbf{x}_C)$ to lower-dimensional manifolds. Parameterized circuits exploit this through an alternating-layer structure that implements a tensor-network decomposition~\cite{huggins2019towards}: single-qubit rotations parameterize local marginals, while entangling gates capture correlations---effectively factorizing the distribution in a manner analogous to how MRFs factorize over cliques. When entanglement matches graphical structure (the clique strategy), this factorization is particularly efficient.
\end{itemize}

\subsection{Strengths of Our Approach}
\begin{itemize}[leftmargin=*, topsep=0pt, partopsep=0pt, itemsep=2pt, parsep=0pt]
    \item \textbf{Research-ready implementation.} Included benchmark runners and artifact checks validate training behavior, performance trends, and end-to-end functionality across representative model structures.
    \item \textbf{Dual-mode flexibility.} The codebase supports exact-style amplitude state preparation for small-$n$ settings and variational compression studies for larger-$n$ settings, enabling accuracy/scalability trade-off analysis under a unified experiment pipeline.
    \item \textbf{Performance optimizations.} Sparse-diagonal construction and careful backend handling make the current simulator-based workflow practical on commodity hardware for the published artifact regimes.
    \item \textbf{Research integrity.} Honest benchmarking distinguishes scenarios where quantum and classical methods are equivalent from those where quantum-specific advantages dominate, avoiding inflated or misleading claims.
    \item \textbf{Reproducibility.} Fixed random seeds, documented hyperparameters, and reproducible example scripts facilitate independent verification and extension.
\end{itemize}

\subsection{Limitations and Failure Modes}\label{sec:limitations_mcmc}
\begin{itemize}[leftmargin=*, topsep=0pt, partopsep=0pt, itemsep=2pt, parsep=0pt]
    \item \textbf{Scope of the studied regime.} The amplitude-encoding pipeline requires $O(2^n)$ classical preprocessing and is restricted to small-$n$ settings where this cost is acceptable---the same regime in which exact inverse-CDF sampling is feasible. We make no claims about quantum performance outside this regime; in particular, we do not implement the QCGM construction~\cite{piatkowski2024qcgms} that could achieve an exponential speedup via Hamiltonian simulation, nor do we evaluate on noisy hardware.
    \item \textbf{Fidelity decay for large models.} Fidelity generally declines as $n$ increases (see Table~\ref{tab:mps_scaling}), highlighting the expressivity/optimization limits of shallow ans\"atze.
    \item \textbf{Optimization scaling limits.} Parameter-shift training requires repeated statevector evaluations; for larger regimes, improved optimizers and adaptive learning-rate schedules remain important future work.
    \item \textbf{No physical-hardware noise modeling.} Our experiments use simulator backends (Aer and BlueQubit simulator services). Deployment on real quantum hardware will require readout correction, zero-noise extrapolation, and dynamical decoupling.
    \item \textbf{What the ESS comparison does and does not show.} The ESS experiments quantify autocorrelation costs for four classical baselines under this artifact protocol. The parallel-tempering result ($1.79\times$ mean ratio) is the most conservative and appropriate headline. The amortized wall-clock comparison against exact inverse-CDF (Table~\ref{tab:ess_per_sec}) further shows that any structural advantage of the quantum sampler does not translate into raw throughput in the regime where the full distribution can be precomputed. These results define the scope of what this submission can and cannot demonstrate; they do not rule out a different conclusion at larger $n$ where $O(2^n)$ preprocessing is infeasible.
\end{itemize}

\subsection{Method Selection Guidelines}
\begin{itemize}[leftmargin=*, topsep=0pt, partopsep=0pt, itemsep=2pt, parsep=0pt]
    \item \textbf{Use exact inverse-CDF sampling} when $n$ is small enough that $O(2^n)$ preprocessing is feasible. This is the dominant classical choice on raw wall-clock in this regime once preprocessing is amortized.
    \item \textbf{Use amplitude-encoded quantum sampling} in the same $n$-regime when (a) downstream workflows specifically require i.i.d.\ samples through a quantum measurement (e.g., quantum-statistical testing), (b) the preprocessing cost can be amortized across many sampling tasks, or (c) integration with a larger quantum pipeline is desired. Do not select it on wall-clock grounds against inverse-CDF.
    \item \textbf{Use parallel tempering or block-structured MCMC} when $n$ is too large for full enumeration but the distribution is peaked enough that single-site Gibbs mixes slowly.
    \item \textbf{Use classical junction-tree or message-passing methods} when the graph is tree-structured or has small treewidth~\cite{lauritzen1988local}.
    \item \textbf{Variational quantum compression (depth-3 hardware-efficient ansatz)} is not a recommended sampling method based on our results: at every $n \le 12$ tested, mean-field or MPS produces better full-distributional fidelity at matched or smaller parameter budgets.
\end{itemize}

\subsection{Future Directions}
\begin{itemize}[leftmargin=*, topsep=0pt, partopsep=0pt, itemsep=2pt, parsep=0pt]
    \item \textbf{Full QCGM implementation.} Ancilla-based circuits~\cite{piatkowski2024qcgms} could enable a true exponential speedup via Hamiltonian simulation ($2$--$4$ weeks estimated effort).
    \item \textbf{Hardware deployment.} Running on devices such as IBM Quantum and IonQ will require readout correction, zero-noise extrapolation, and dynamical decoupling. Experiments on $50$--$100$-qubit devices would provide the first practical performance data for quantum MRF sampling.
    \item \textbf{Marginal-matching techniques.} Enforcing clique marginals $\hat{P}(\mathbf{x}_C) \approx P_\theta(\mathbf{x}_C)$ during variational training could improve fidelity for large models by an estimated $0.1$--$0.15$.
    \item \textbf{Hybrid quantum-classical sampling.} Using quantum circuits for high-probability regions and classical MCMC for tail exploration could reduce computational cost while maintaining sample quality.
    \item \textbf{Advanced optimization strategies.} For $n > 15$, lower-variance gradient estimators, improved optimizers, and schedule tuning are promising paths to more efficient training.
\end{itemize}
Longer-term, implementation of full QCGM on 50--100-qubit devices with error mitigation could bridge the gap between quantum simulation and practical quantum applications for probabilistic inference in computer vision, computational biology, and generative modeling.

\section{Conclusion}

We presented a characterization of autocorrelation costs for four classical sampling baselines on discrete MRFs, compared against an amplitude-encoded quantum sampler that produces structurally i.i.d.\ output. Across 60 instances and five graph families, the Quantum/baseline ESS ratios decrease monotonically from $16.35$ (single-site Gibbs) to $7.29$ (block Gibbs) to $1.82$ (tuned-block Gibbs) to $1.79$ (parallel tempering), making clear how much of the apparent quantum sampling advantage in the literature is attributable to baseline choice. Once the $O(2^n)$ classical preprocessing required by the amplitude encoder is amortized into the wall-clock comparison, exact inverse-CDF sampling achieves $153\times$ the amortized ESS/s of the quantum sampler. Together these results define a precise scope within which the amplitude encoder is a defensible sampling tool.

Secondary results: the amplitude-encoding pipeline reaches $F\approx1.000$ at $n=8,10,12$ on simulator backends (a sanity check on a Qiskit primitive); a multi-trial MPS scaling study through $n=40$ establishes a classical ceiling of $F=0.721\pm0.059$ at $\chi=32$; and a shallow hardware-efficient VQC underperforms matched-budget MPS at every $n\le12$ tested. Several diagnostic null results are reported and we believe are useful to the community.

We do not claim quantum advantage in this work. The most plausible paths to one---full QCGM Hamiltonian simulation, deeper or problem-aware variational ans\"atze, or noisy-hardware benchmarks at scale---are clearly identified as future work. All code, data, and artifacts are publicly available at \url{https://github.com/arulrhikm/QuantumDGM} under an MIT license.

\section*{Reproducibility Note}
This manuscript inlines artifact values from \texttt{paper\_table\_metrics.json} generated on \texttt{2026-04-27 09:03:28 UTC}. The file lives at the project root and at \texttt{figures/paper\_table\_metrics.json} (generated copy). Experiment seeds are in \texttt{experiments/seeds.json}. To regenerate, run:
\begin{verbatim}
python -m experiments.run_all_for_paper --fast
python -m experiments.generate_paper_artifacts
\end{verbatim}
then copy updated values into this self-contained \texttt{paper.tex}.

\section*{Acknowledgments}

We thank the open-source quantum computing community for the tools that enabled this work. Code and documentation are available under an MIT license at \url{https://github.com/arulrhikm/QuantumDGM}.

\begin{thebibliography}{00}

\bibitem{li2009markov} S.~Z. Li, \emph{Markov Random Field Modeling in Image Analysis}, 3rd ed. London, U.K.: Springer, 2009.

\bibitem{kamisetty2013coevolution} H.~Kamisetty, S.~Ovchinnikov, and D.~Baker, ``Assessing the utility of coevolution-based residue-residue contact predictions in a sequence- and structure-rich era,'' \emph{Proc. Natl. Acad. Sci. U.S.A.}, vol.~110, no.~39, pp.~15674--15679, 2013.

\bibitem{salakhutdinov2009deep} R.~Salakhutdinov and G.~Hinton, ``Deep Boltzmann machines,'' in \emph{Proc. AISTATS}, 2009, pp.~448--455.

\bibitem{geman1984stochastic} S.~Geman and D.~Geman, ``Stochastic relaxation, Gibbs distributions, and the Bayesian restoration of images,'' \emph{IEEE Trans. Pattern Anal. Mach. Intell.}, vol.~PAMI-6, no.~6, pp.~721--741, 1984.

\bibitem{wainwright2008graphical} M.~J. Wainwright and M.~I. Jordan, ``Graphical models, exponential families, and variational inference,'' \emph{Found. Trends Mach. Learn.}, vol.~1, no.~1--2, pp.~1--305, 2008.

\bibitem{pearl1988probabilistic} J.~Pearl, \emph{Probabilistic Reasoning in Intelligent Systems}. San Francisco, CA: Morgan Kaufmann, 1988.

\bibitem{piatkowski2024qcgms} N.~Piatkowski and C.~Zoufal, ``On quantum circuits for discrete graphical models,'' \emph{Quantum Machine Intelligence}, vol.~6, art.~37, 2024.

\bibitem{amin2018qbm} M.~H. Amin, E.~Andriyash, J.~Rolfe, B.~Kulchytskyy, and R.~Melko, ``Quantum Boltzmann machine,'' \emph{Phys. Rev. X}, vol.~8, no.~2, p.~021050, 2018.

\bibitem{low2014bayes} G.~H. Low, T.~J. Yoder, and I.~L. Chuang, ``Quantum inference on Bayesian networks,'' \emph{Phys. Rev. A}, vol.~89, no.~6, p.~062315, 2014.

\bibitem{verdon2019vqe} G.~Verdon, J.~Marks, S.~Nanda, S.~Leichenauer, and J.~Hidary, ``Quantum Hamiltonian-based models and the variational quantum thermalizer algorithm,'' \emph{arXiv:1910.02071}, 2019.

\bibitem{layden2023qemcmc} D.~Layden, G.~Mazzola, R.~V. Mishmash, M.~Motta,
P.~Wocjan, J.-S. Kim, and S.~Sheldon, ``Quantum-enhanced Markov chain Monte
Carlo,'' \emph{Nature}, vol.~619, pp.~282--287, 2023.

\bibitem{kandala2017hardware} A.~Kandala \emph{et al.}, ``Hardware-efficient variational quantum eigensolver for small molecules and quantum magnets,'' \emph{Nature}, vol.~549, pp.~242--246, 2017.

\bibitem{schuld2019evaluating} M.~Schuld, V.~Bergholm, C.~Gogolin, J.~Izaac, and N.~Killoran, ``Evaluating analytic gradients on quantum hardware,'' \emph{Phys. Rev. A}, vol.~99, no.~3, p.~032331, 2019.

\bibitem{mcclean2018barren} J.~R. McClean, S.~Boixo, V.~N. Smelyanskiy, R.~Babbush, and H.~Neven, ``Barren plateaus in quantum neural network training landscapes,'' \emph{Nat. Commun.}, vol.~9, no.~1, p.~4812, 2018.

\bibitem{huggins2019towards} W.~Huggins, P.~Patel, K.~B. Whaley, and E.~M. Stoudenmire, ``Towards quantum machine learning with tensor networks,'' \emph{Quantum Sci. Technol.}, vol.~4, no.~2, p.~024001, 2019.

\bibitem{liu2018differentiable} J.-G. Liu and L.~Wang, ``Differentiable learning of quantum circuit Born machines,'' \emph{Phys. Rev. A}, vol.~98, p.~062324, 2018.

\bibitem{benedetti2019generative} M.~Benedetti, E.~Lloyd, S.~Sack, and M.~Fiorentini, ``A generative modeling approach for benchmarking and training shallow quantum circuits,'' \emph{npj Quantum Information}, vol.~5, art.~45, 2019.

\bibitem{wittek2017quantum} P.~Wittek and C.~Gogolin, ``Quantum enhanced inference in Markov logic networks,'' \emph{Scientific Reports}, vol.~7, art.~45672, 2017.

\bibitem{ferguson2025qemcmc} R.~Ferguson and P.~Wallden, ``Quantum-enhanced MCMC for systems larger than your quantum computer,'' \emph{Phys. Rev. Research}, vol.~7, p.~013231, 2025.

\bibitem{geyer1992practical} C.~J. Geyer, ``Practical Markov chain Monte Carlo,'' \emph{Stat. Sci.}, vol.~7, no.~4, pp.~473--483, 1992.

\bibitem{brooks2011handbook} S.~Brooks, A.~Gelman, G.~L. Jones, and X.-L. Meng, Eds., \emph{Handbook of Markov Chain Monte Carlo}. Boca Raton, FL: CRC Press, 2011.

\bibitem{lauritzen1988local} S.~L. Lauritzen and D.~J. Spiegelhalter, ``Local computations with probabilities on graphical structures and their application to expert systems,'' \emph{J. R. Stat. Soc. Ser. B}, vol.~50, no.~2, pp.~157--224, 1988.

\end{thebibliography}

\end{document}